\documentclass[11pt]{article}

% ---------- packages ----------
\usepackage[utf8]{inputenc}
\usepackage[T1]{fontenc}
\usepackage{lmodern}
\usepackage[a4paper,margin=1in]{geometry}
\usepackage{amsmath,amssymb}
\usepackage{booktabs}
\usepackage{array}
\usepackage{tabularx}
\usepackage{microtype}
\usepackage{enumitem}
\usepackage{xcolor}
\usepackage{graphicx}
\usepackage{titlesec}
\usepackage[hidelinks,colorlinks=true,linkcolor=teal!60!black,citecolor=teal!60!black,urlcolor=teal!60!black]{hyperref}

% ---------- light styling ----------
\definecolor{accent}{RGB}{20,90,90}
\titleformat{\section}{\large\bfseries\color{accent}}{\thesection}{0.6em}{}
\titleformat{\subsection}{\normalsize\bfseries}{\thesubsection}{0.6em}{}
\setlist{itemsep=2pt,topsep=3pt,parsep=0pt}
\renewcommand{\arraystretch}{1.25}
\newcolumntype{Y}{>{\raggedright\arraybackslash}X}

\title{\vspace{-1.2cm}\textbf{Domain Decomposition for Neural Surrogate Models}\\[2pt]
\large A survey of approaches for predicting large fields from shared subdomain networks}
\author{Research note --- \texttt{pyurbanair} / domain-decomposition surrogates}
\date{\today}

\begin{document}
\maketitle
\vspace{-0.6cm}

\begin{abstract}
\noindent
This note surveys methods for building neural surrogate models that operate on
\emph{subdomains} of a large simulation rather than the full field at once, and
that recombine subdomain predictions into a global solution. The motivation is
practical: a single-shot surrogate over a full three-dimensional large-eddy
simulation (LES) field is infeasible at the scales of interest, whereas a
network trained to act on smaller patches can be reused across an arbitrarily
large domain. The literature relevant to this idea is spread across numerical
analysis (Schwarz and substructuring methods), scientific machine learning
(domain-decomposed PINNs), operator learning (Fourier and graph neural
operators), computer vision (overlap-tile inference), data-driven weather
prediction (limited-area models with boundary forcing), and ML for turbulence.
The note organizes these into a small set of design axes --- how the domain is
split, how interface information is exchanged, and how subdomain outputs are
merged --- and closes with the specific implications for the \texttt{pyurbanair}
surrogate, which is a translation-invariant 3D convolutional one-step predictor,
and a recommended staged approach for it. It is intended as an overview to
support a decision on which approach to pursue.
\end{abstract}

\section{Motivation and problem setting}

The \texttt{pyurbanair} neural surrogate is a learned one-step time-stepper for
urban CFD. A convolutional network (a 3D UNet with ConvNeXt blocks, or a simpler
convolutional baseline) maps a state $\mathbf{s}_n$ on a regular cell-centered
grid, conditioned on inflow parameters and a binary geometry mask, to the next
state $\mathbf{s}_{n+1}$; full trajectories are produced by autoregressive
rollout. Training uses one-step transition pairs with an optional pushforward
unroll for rollout stability \cite{brandstetter2022mppde}. At present the
network models a single, fixed domain defined by the training data: the trained
grid resolution and bounds are checked at inference and a mismatch is rejected.

This whole-domain design does not scale. A 3D LES field of a realistic urban
area has too many cells to fit a useful receptive field and the activations of a
deep 3D network into memory, and the cost of a single forward pass grows with
the volume. The proposal is to train \emph{one} network on subdomains (patches)
of the field and apply it to many subdomains, combining the local predictions
into a global state. The convolutional surrogate is well suited to this because
convolutions are translation-equivariant: the same kernels are valid at any
spatial location, so a patch-trained network can in principle be applied to any
patch of a larger field without retraining. What is not automatic is (i) how
each subdomain obtains the boundary information it needs from its neighbors, and
(ii) how overlapping or abutting predictions are reconciled into a single
consistent field, especially across many autoregressive steps.

Those two questions --- interface coupling and merging --- are exactly the
questions that classical domain decomposition has studied for decades, and that
several recent machine-learning communities have re-encountered independently.
The rest of this note reviews how each field answers them.

\section{Design axes}\label{sec:axes}

Across very different methods, the same choices recur. It is useful to fix them
up front; the taxonomy table in Section~\ref{sec:synthesis} places the surveyed
methods along these axes.

\begin{description}[leftmargin=1.2em,style=nextline]
\item[Partition geometry.] \emph{Overlapping} subdomains share a strip of cells
with their neighbors; \emph{non-overlapping} subdomains meet at a shared
interface of measure zero. Overlap costs redundant computation but simplifies
continuity; non-overlapping partitions need explicit interface conditions.
\item[Interface information.] How a subdomain learns about its neighbors. Options
seen in the literature: a \emph{halo} (ghost cells / boundary strip) supplied as
input; \emph{interface loss terms} penalizing discontinuity; \emph{Schwarz
iteration} exchanging boundary traces until convergence; \emph{message passing}
along a graph; or a built-in \emph{partition-of-unity} ansatz that is continuous
by construction.
\item[Merging.] How local outputs become one field. Options: \emph{exact
center-cropping} (keep only the interior of each tile, valid when the halo is at
least a receptive field wide); \emph{overlap averaging / partition-of-unity
blending} with a smooth window; \emph{Schwarz assembly} after iteration; or
provably \emph{consistent} distributed evaluation where the tiled result equals
the whole-domain result.
\item[Network sharing.] One shared network reused across all subdomains, versus
one network per subdomain. Sharing is the norm in operator-learning and CV
tiling; per-subdomain networks are common in PINN decompositions, where each
network fits one local solution.
\item[Coupling regime.] \emph{One-shot} (each subdomain predicted once, then
merged) versus \emph{iterative} (subdomain solves repeated until interfaces
agree). Iterative methods are accurate but expensive; one-shot methods are cheap
but can leave seams.
\end{description}

\section{Classical domain decomposition}

Domain decomposition methods (DDM) were developed to solve PDE discretizations
in parallel by splitting the domain into subproblems. They provide both the
vocabulary and the convergence theory that the ML methods reuse, and the
standard references are \cite{toselli2005ddm,quarteroni1999ddm,dolean2015ddm}.

Two families matter here. \emph{Overlapping Schwarz methods} go back to the
alternating scheme analyzed by Lions \cite{lions1988schwarz}: each subdomain is
solved with its neighbors' current solution imposed as Dirichlet data on the
internal interface, and the process is iterated. The additive (parallel) variant
with a coarse correction \cite{dryja1994smalloverlap} is the workhorse
preconditioner; convergence improves with overlap width and, crucially, depends
on a coarse space to avoid degrading as the number of subdomains grows.
\emph{Non-overlapping substructuring methods} instead tear the domain at shared
interfaces and reconnect it: FETI \cite{farhat1991feti} and its dual--primal
successor FETI-DP \cite{farhat2001fetidp} enforce interface continuity with
Lagrange multipliers, while the balancing methods BDD \cite{mandel1993bdd} and
BDDC \cite{dohrmann2003bddc} use primal coarse constraints on corners, edges and
faces. The recurring lesson, independent of the variant, is that purely local
coupling does not propagate information globally fast enough: a coarse space (a
second, global level) is needed for the iteration count to stay bounded as the
domain is cut into more pieces. This reappears below in multilevel PINNs and in
multipole graph operators.

A further idea that the ML work has adopted is the \emph{optimized Schwarz}
transmission condition: replacing Dirichlet interface data with a Robin
combination of value and flux accelerates convergence, and the Robin parameter
can be tuned to the problem \cite{dolean2015ddm}.

\section{Domain decomposition in physics-informed neural networks}

PINNs \cite{raissi2019pinn} fit a network to satisfy a PDE residual plus
boundary data. A single global network struggles on large or multi-scale
domains, which motivated decomposed variants; the review by Heinlein et al.
\cite{heinlein2021review} is the standard entry point and separates two
directions: using ML inside classical DDM (for example, predicting adaptive
coarse-space constraints \cite{heinlein2019adaptive}), and using DDM to scale ML
solvers. The second direction is the relevant one.

\emph{Non-overlapping, interface-loss methods.} Conservative PINNs (cPINN)
\cite{jagtap2020cpinn} place a separate network on each non-overlapping
subdomain and stitch them with interface terms enforcing flux continuity and
solution-average matching, preserving conservation for conservation laws.
Extended PINNs (XPINN) \cite{jagtap2020xpinn} generalize this to arbitrary
space--time partitions and any PDE, with residual and average-continuity
interface losses. Both allow per-subdomain network sizing and parallel training;
the MPI-parallel implementation is described in
\cite{shukla2021parallelpinn}. The cost is that continuity is only enforced
weakly, through penalty terms whose weighting must be tuned.

\emph{Overlapping, partition-of-unity methods.} Finite basis PINNs (FBPINN)
\cite{moseley2023fbpinn} take the opposite route: subdomains overlap, and each
local network is multiplied by a smooth, compactly supported window function so
that the global solution is a partition-of-unity sum of local networks. This is
continuous by construction --- no interface loss is needed --- and per-subdomain
input normalization mitigates spectral bias. The single-level version still
suffers the classical scalability limit, which the multilevel extension
\cite{dolean2024multilevel} addresses by stacking nested overlapping
decompositions, exactly mirroring multilevel Schwarz.

\emph{Schwarz-iterative methods.} DeepDDM \cite{li2020deepddm} replaces the
subdomain solvers in an overlapping Schwarz iteration with PINNs: each subdomain
is solved by a network, interface traces are exchanged, and the outer iteration
is run to convergence; it inherits the classical behavior (iteration count
roughly independent of network size, improving with overlap). The two-level
extension \cite{dolean2024twolevel} adds a coarse network to recover scalability.
The variational D3M \cite{li2020d3m} couples local networks through an energy
formulation. Basir and Senocak \cite{basir2023schwarz} use a non-overlapping
generalized Schwarz method with \emph{learnable} Robin interface parameters, and
Snyder et al.\ \cite{snyder2023schwarzpinn} use the Schwarz alternating method to
couple PINNs both to one another and to conventional FEM/FD solvers --- a
template for hybrid neural/numerical coupling. Augmented PINNs (APINN)
\cite{hu2023apinn} soften the hard partition into a trainable gating network, a
learned partition of unity with partial weight sharing.

For a surrogate (which learns a map, not a single solution) these methods are
mostly relevant for their interface treatments rather than their training
setups: PINN decompositions typically fit one network per subdomain to one
problem instance, whereas the goal here is one shared network that generalizes.

\section{Neural operators and the locality question}

Operator-learning surrogates map function to function and, unlike PINNs,
generalize across problem instances --- the regime of interest here. The Fourier
Neural Operator (FNO) \cite{li2021fno} is the prominent example, but its spectral
convolution is \emph{global}: truncated Fourier modes mix information across the
entire domain, and the architecture assumes a single rectangular, periodic,
uniform grid. This is precisely the property that obstructs tiling. The unifying
theory of neural operators \cite{kovachki2023neuraloperator} establishes that one
parameter set defines a discretization-invariant operator, which is the formal
justification for reusing a single network across subdomains of differing size.

Several lines address locality. Factorized FNO \cite{tran2023ffno} makes the
spectral layers separable, cutting memory enough to reach 3D; Geo-FNO
\cite{li2023geofno} learns a deformation to handle non-rectangular geometries;
and localized-kernel operators \cite{liuschiaffini2024localized} add explicit
local stencils to recover fine structure the global spectral layer misses. The
local neural operator \cite{ye2024lno} is shift-invariant with a finite
receptive field, so it applies to domains it was not trained on; its companion
analysis quantifies how large that receptive field --- effectively, how large a
halo --- must be to avoid spurious oscillation, which is directly useful for
sizing overlaps.

Three genuinely different meanings of ``decomposition'' appear in the operator
literature and should not be conflated. \emph{Model-parallel} FNO
\cite{grady2023modelparallelfno} shards a single global operator's data and
weights across GPUs to reach billions of variables; because the FFT is global,
each spectral layer requires halo/all-to-all communication, and the merged
output is exact rather than approximate. This is parallelism for scale, not
approximate stitching. \emph{Schwarz-type operator coupling} is the closer
analogue: Huang et al.\ \cite{huang2025operatorddm} partition the domain, solve
local problems with a shared neural operator, and stitch them with an iterative
Schwarz neural inference scheme, with convergence and error analysis and strong
geometry generalization; the structural-analysis DLD3 framework
\cite{klawonn2024dld3} similarly wraps a shared local network in an overlapping
Schwarz iteration.

\section{Graph and message-passing surrogates: halo exchange}

Graph-based simulators make locality explicit and are the cleanest setting for
halo handling. Graph network simulators \cite{sanchezgonzalez2020gns} and
MeshGraphNets \cite{pfaff2021meshgraphnets} run message passing over a particle
or mesh graph, so each node depends only on a bounded neighborhood; message
passing neural PDE solvers \cite{brandstetter2022mppde} add the pushforward trick
for autoregressive stability (the same trick used in \texttt{pyurbanair}). Graph
kernel \cite{li2020gno} and multipole graph \cite{li2020mgno} operators frame
this as operator learning, with the multipole version using a coarsening
hierarchy to capture long-range coupling at linear cost --- the operator-learning
analogue of a coarse space.

For scaling, the relevant references partition a large mesh and process subgraphs
with a shared network. X-MeshGraphNet \cite{nabian2024xmeshgraphnet} augments each
partition with a halo region so subgraphs can be processed independently, the
halo supplying the boundary information that would otherwise cross the partition.
The distributed GNN of Barwey et al.\ \cite{barwey2024consistentgnn} goes further
with a consistent message-passing layer using halo-node exchange, so that
evaluation on many partitions is arithmetically identical to evaluation on the
whole graph --- a provably seam-free merge, demonstrated against a spectral-element
CFD solver. CoMLSim-style latent decomposition \cite{hanna2024comlsim} encodes
each subdomain into a latent vector and exchanges latent information between
neighbors to enforce coupling, time-stepping autoregressively in latent space;
it is the closest published match to ``a shared network tiled over subdomains,
stepping in time.''

\section{Tiling mechanics from computer vision}

Computer vision has applied a single network to images larger than memory for a
decade, and its stitching mechanics transfer directly. The overlap-tile strategy
introduced with U-Net \cite{ronneberger2015unet} predicts each output tile from a
larger input tile (the surrounding context is the halo); with unpadded
convolutions the output tiles abut exactly, so no blending is needed, and the
image border is handled by mirror padding. When exactness is required, keeping
only the receptive-field-valid interior of each tile reproduces whole-image
inference bit-for-bit. Where tiles do overlap, predictions are blended with a
smooth partition-of-unity window (a Hann or Gaussian taper, weight one at the
center decaying to zero at the edge); generative pipelines such as MultiDiffusion
\cite{bartal2023multidiffusion} average overlapping tile predictions at every
denoising step to produce seamless panoramas from one shared model. A subtle but
citable pitfall: tiling artifacts often originate not in the overlap policy but
in per-tile normalization statistics that differ between tiles
\cite{buglakova2025tiling}; a convolutional surrogate with normalization layers
can produce seams for the same reason. Fully convolutional super-resolution GANs
trained on patches generalize to large fields at inference
\cite{stengel2020phiregan}, evidence that patch-trained generators tile if their
statistics are controlled.

\section{Limited-area weather models and boundary forcing}

Data-driven weather prediction faces the same scaling question on the sphere and
has converged on instructive answers. Global models --- FourCastNet
\cite{pathak2022fourcastnet} (adaptive FNO), GraphCast \cite{lam2023graphcast}
(multi-mesh GNN), Pangu-Weather \cite{bi2023pangu} (3D windowed-attention
transformer) --- process the whole field at once and handle locality implicitly,
serving as the ``no decomposition'' baseline. The directly relevant work is the
\emph{limited-area model} (LAM): Neural-LAM \cite{oskarsson2023neurallam}
forecasts only an interior region and supplies lateral boundary conditions as
forcing along a boundary strip at each autoregressive step --- the ML counterpart
of a halo or relaxation zone fed from an external model. The MLLAM framework
\cite{oskarsson2025mllam} and the operational AI-LAM of Nipen et al.
\cite{nipen2025ailam} generalize boundary forcing, the latter taking its lateral
boundaries from a global AI model in a one-way nesting. An alternative avoids
explicit boundaries entirely: a global \emph{stretched grid}
\cite{husain2024stretchedgrid} refines resolution over the region of interest so
that systems pass seamlessly between coarse and fine zones within one continuous
field. For \texttt{pyurbanair}, the LAM boundary-forcing pattern is the most
transferable: it shows how to feed halo/interface state into an autoregressive
surrogate each step rather than only at initialization.

\section{Turbulence- and LES-specific surrogates}

Within fluids, several works bear directly on tiling a turbulence network. The
most on-point is the parallel super-resolution GAN coupled to an LES solver for
subgrid reconstruction, which studies the halo (overlap) width required for
accurate reconstruction at subdomain boundaries and reports scaling as a function
of subdomain size; it finds training sub-boxes must be large relative to the
Kolmogorov scale, a concrete constraint on minimum patch size. The convolutional
super-resolution of turbulence by Fukami et al.\
\cite{fukami2019superresolution,fukami2021spatiotemporal} establishes that
fully-convolutional networks reconstruct turbulent fields from coarse input and
can be applied locally; learned coarse turbulence models with dilated-convolution
processors \cite{stachenfeld2022coarse} are translation-equivariant with a
bounded receptive field, the same property \texttt{pyurbanair} relies on.
Multi-agent RL wall models \cite{bae2022scimarl} decompose the near-wall region
into local agents sharing one policy --- a subdomain structure with a single
shared network. Finally, the difficulty of long autoregressive rollouts in 3D
turbulence motivates generative alternatives such as TurbDiff
\cite{lienen2024turbdiff} and autoregressive conditional diffusion
\cite{kohl2023benchmarking}, and transformer neural operators are now being
applied to LES closure directly \cite{li2024tno}. These are relevant because
tiling in space compounds with stepping in time: errors can accumulate both at
interfaces and over the rollout.

\section{Synthesis}\label{sec:synthesis}

Table~\ref{tab:taxonomy} places representative methods on the axes of
Section~\ref{sec:axes}. Table~\ref{tab:merge} summarizes the merge strategies,
which is the design decision most specific to the surrogate setting.

\begin{table}[t]
\centering\footnotesize
\setlength{\tabcolsep}{4pt}
\begin{tabularx}{\textwidth}{@{}l l l l l@{}}
\toprule
\textbf{Method} & \textbf{Partition} & \textbf{Interface} & \textbf{Merge} & \textbf{Network} \\
\midrule
Overlapping Schwarz \cite{lions1988schwarz} & overlap & Dirichlet trace, iterate & assembly & --- \\
FETI-DP / BDDC \cite{farhat2001fetidp,dohrmann2003bddc} & none & multipliers + coarse & exact & --- \\
cPINN / XPINN \cite{jagtap2020cpinn,jagtap2020xpinn} & none & residual + avg.\ loss & weak (loss) & per-subdomain \\
FBPINN \cite{moseley2023fbpinn} & overlap & partition-of-unity ansatz & PoU sum & per-subdomain \\
Multilevel FBPINN \cite{dolean2024multilevel} & overlap, nested & PoU + coarse levels & PoU sum & per-sub./level \\
DeepDDM \cite{li2020deepddm} & overlap & Schwarz trace, iterate & assembly & per-subdomain \\
Operator DDM (SNI) \cite{huang2025operatorddm} & overlap & Schwarz, iterate & assembly & \textbf{shared} \\
Model-parallel FNO \cite{grady2023modelparallelfno} & sharded & global halo (FFT) & exact & shared (1 op) \\
X-MeshGraphNet \cite{nabian2024xmeshgraphnet} & none & halo region & (consistent) & \textbf{shared} \\
Consistent dist.\ GNN \cite{barwey2024consistentgnn} & none & halo-node exchange & exact & \textbf{shared} \\
CoMLSim latent DDM \cite{hanna2024comlsim} & none & latent exchange & assembly & \textbf{shared} \\
U-Net overlap-tile \cite{ronneberger2015unet} & overlap (input) & input halo & center-crop & \textbf{shared} \\
MultiDiffusion \cite{bartal2023multidiffusion} & overlap & --- & PoU average & \textbf{shared} \\
Neural-LAM \cite{oskarsson2023neurallam} & region + halo & boundary forcing & --- (single region) & \textbf{shared} \\
\bottomrule
\end{tabularx}
\caption{Representative methods on the design axes of Section~\ref{sec:axes}.
``Shared'' (bold) marks a single network reused across subdomains --- the regime
relevant to \texttt{pyurbanair}.}
\label{tab:taxonomy}
\end{table}

\begin{table}[t]
\centering\small
\begin{tabularx}{\textwidth}{@{}l Y Y@{}}
\toprule
\textbf{Merge strategy} & \textbf{How it works} & \textbf{Trade-off} \\
\midrule
Exact center-crop & Overlap input by $\geq$ one receptive field; keep only each tile's valid interior. & Seam-free and cheap; needs known, finite receptive field and wide halos. \\
PoU / overlap averaging & Weight tiles by a smooth window (Hann/Gaussian) and sum in overlaps. & Simple, one-shot, robust to small mismatches; small blurring in overlaps. \\
Schwarz assembly & Iterate subdomain solves exchanging interface data until convergence. & Accurate, theory available; iterative cost, harder for transient rollout. \\
Consistent distributed & Halo-node exchange makes tiled evaluation equal whole-domain. & Provably seam-free; requires architecture and halo bookkeeping. \\
Boundary forcing & Predict interior only; feed neighbor/coarse state on a boundary strip each step. & Natural for autoregressive stepping; needs a boundary source. \\
\bottomrule
\end{tabularx}
\caption{Merge strategies for recombining subdomain predictions.}
\label{tab:merge}
\end{table}

Three cross-cutting points stand out. First, a single shared network across
subdomains is standard in every operator-learning, graph, and CV tiling method,
and is well-justified by translation invariance and discretization-invariance
theory \cite{kovachki2023neuraloperator}; the per-subdomain networks of the PINN
literature are an artifact of fitting one solution, not a requirement.
Second, every field rediscovers that purely local coupling under-propagates
information, and adds either a coarse level
\cite{dryja1994smalloverlap,dolean2024multilevel,li2020mgno} or wide enough
halos \cite{ye2024lno,barwey2024consistentgnn} to compensate. Third, the merge
strategy and the coupling regime (one-shot vs.\ iterative) are the levers with
the largest effect on cost and seam quality.

\section{Implications for \texttt{pyurbanair}}

The current surrogate has properties that make several of these approaches
directly applicable. It is convolutional and therefore translation-equivariant,
so a single network can be tiled without retraining; it already conditions on a
per-cell geometry mask, which means subdomain-specific urban geometry enters
naturally as an input channel; and it already uses the pushforward trick for
rollout stability. The main open design choices map onto the axes above.

\paragraph{A pragmatic starting point.} Train on patches with a halo, and at
inference tile the domain with overlapping patches, predict each with the shared
network, and merge by either exact center-cropping (if the effective receptive
field is known and the halo covers it) or partition-of-unity blending in the
overlaps. This is the U-Net overlap-tile recipe \cite{ronneberger2015unet} plus
MultiDiffusion-style averaging \cite{bartal2023multidiffusion}, and it requires
no architectural change. The receptive-field analysis of local neural operators
\cite{ye2024lno} and the LES halo-width study give concrete guidance on the
minimum overlap.

\paragraph{Feeding the halo each step.} Tiling in space interacts with stepping
in time: a one-step network needs valid neighbor state at the patch boundary at
\emph{every} step, not only at initialization. Two options exist. Supply the halo
as a boundary-forcing input strip taken from neighboring patches' previous-step
predictions, in the Neural-LAM style \cite{oskarsson2023neurallam}; or adopt an
explicitly halo-aware evaluation with neighbor exchange between steps, as in the
consistent distributed GNN \cite{barwey2024consistentgnn}. The former is simpler
and one-shot; the latter is closer to seam-free but needs exchange bookkeeping.

\paragraph{The global-coupling caveat.} The most important physics caveat is that
incompressible (and low-Mach) flow has an elliptic pressure component: the
pressure field couples the whole domain instantaneously, so a purely local
tile-by-tile predictor will miss long-range coupling that the underlying LES
resolves. Every numerical and ML domain decomposition that targets elliptic
problems addresses this with a coarse space or a global level
\cite{dryja1994smalloverlap,dolean2024multilevel,li2020mgno}. A serious
\texttt{pyurbanair} decomposition will likely need an analogous mechanism --- a
coarse global network, a multipole/multi-scale path, or a Schwarz iteration on
the pressure --- rather than independent local rollouts alone. This is the single
point most worth resolving before committing to an approach.

\paragraph{Geometry and positional conditioning.} Because urban geometry varies
between subdomains, the geometry mask channel is essential, and it may help to add
positional or boundary-distance encodings so the network can distinguish an
interior patch from a domain-edge patch. The CoMLSim latent decomposition
\cite{hanna2024comlsim} and operator DDM \cite{huang2025operatorddm} are the
closest published templates for a shared network stepping over coupled
subdomains, and are worth reading in full before the next decision.

\section{Recommendations}

This section states a recommended path and the reasoning behind it. The reasoning
rests on three facts about the present setup. The surrogate is convolutional and
translation-equivariant, so one network can be tiled without retraining and
per-subdomain networks (the PINN pattern) are unnecessary. The flow is
incompressible, so pressure couples the whole domain through an elliptic
constraint and a purely local predictor will be wrong in a way that does not
shrink with longer training. And the model is a one-step time-stepper rolled out
autoregressively, so a subdomain needs valid neighbor state at its boundary at
every step, not only at initialization.

\subsection*{Recommended target: a two-level overlapping decomposition with a shared local network}

The recommended end state is a two-level scheme. A single shared convolutional
network predicts one step on overlapping fine patches; each patch is conditioned
on a halo strip of neighbor state and on a coarse global context field; patch
predictions are merged by partition-of-unity blending. The coarse level is the
part that distinguishes this from naive tiling, and it is included specifically
to carry the long-range (pressure) coupling that local patches cannot. This
mirrors the coarse-space construction that classical Schwarz
\cite{dryja1994smalloverlap}, multilevel FBPINNs \cite{dolean2024multilevel}, and
multipole graph operators \cite{li2020mgno} all converged on for the same reason.
The local-patch-plus-neighbor-exchange structure follows the latent subdomain
decomposition of CoMLSim \cite{hanna2024comlsim} and the operator domain
decomposition of Huang et al.\ \cite{huang2025operatorddm}, which are the two
closest published templates for a shared network stepping over coupled
subdomains.

Concretely: run the existing network (or a coarsened copy) on a downsampled
full-domain field to produce a cheap, low-resolution global state each step;
interpolate that coarse state onto each fine patch and pass it as additional
input channels alongside the geometry mask and the halo; predict the fine patch;
and blend overlaps with a smooth window. The coarse pass is cheap because it runs
at low resolution, and it gives every patch a consistent view of the global field
without an iterative solve.

\subsection*{A staged path to get there}

Building this in one step is risky, so the recommendation is to stage it and stop
early if a simpler stage is good enough.

\begin{enumerate}[leftmargin=1.4em]
\item \textbf{Patch training and overlap-tile inference (no global coupling).}
Train the current architecture on patches cropped from existing full-domain
simulations, each crop carrying a halo, so the network sees realistic interface
conditions during training rather than only at test time. At inference, tile with
overlap, predict each patch, and merge by partition-of-unity blending
\cite{ronneberger2015unet,bartal2023multidiffusion}. This is the cheapest thing
that could work and requires no architecture change. Use it to measure how far a
purely local predictor gets before the elliptic coupling becomes the dominant
error --- that measurement decides how much of the rest is needed.
\item \textbf{Halo forcing across the rollout.} Feed each patch's boundary strip
from its neighbors' previous-step predictions every step, in the limited-area
weather style \cite{oskarsson2023neurallam,oskarsson2025mllam}. This is a small
change to the rollout loop and addresses the time-stepping interaction that
single-shot tiling ignores.
\item \textbf{Add the coarse global level.} Introduce the low-resolution global
context described above. This is the step expected to matter most for
incompressible urban flow, and the one to reach for as soon as stage~1 shows a
local-coupling error floor.
\item \textbf{Consistent exchange, only if seams persist.} If blended overlaps
still accumulate visible seams over long rollouts, replace blending with explicit
halo-node exchange between patches each step, which can make the tiled evaluation
match a whole-domain evaluation \cite{barwey2024consistentgnn}. This is the most
implementation-heavy option and should not be adopted pre-emptively.
\end{enumerate}

\subsection*{Supporting choices}

A few smaller decisions are worth making deliberately. Predict a residual
(delta-state) rather than the next state directly; the current architectures
predict the state outright, and for a time-stepper a residual target usually
improves rollout stability and is a small change. Add a boundary-distance or
positional encoding as an input channel so the network can tell an interior patch
from a domain-edge patch, since their correct boundary behavior differs. Be
careful with normalization layers: per-patch normalization statistics are a known
and avoidable cause of tiling seams \cite{buglakova2025tiling}, so use
normalization whose statistics do not depend on the patch contents. Size the halo
from the network's effective receptive field and the flow's Kolmogorov scale; the
local-operator receptive-field analysis \cite{ye2024lno} and the LES
super-resolution halo study give the relevant guidance. Evaluate not only
pointwise error but interface continuity, the velocity divergence near patch
boundaries (a direct check on the pressure-coupling problem), and energy spectra
against full-domain references.

\subsection*{What to avoid}

Two tempting options are not recommended as the primary path. A global Fourier
neural operator \cite{li2021fno} is a poor fit despite its popularity for
turbulence: its spectral convolution is global and tied to a single periodic
rectangular grid, which is the opposite of what a tiled urban-geometry surrogate
needs, and the model-parallel route \cite{grady2023modelparallelfno} scales the
whole-domain operator rather than enabling a small shared subdomain network. A
full Schwarz iteration to convergence at every time step
\cite{li2020deepddm,huang2025operatorddm} is accurate but expensive, and the
iteration interacts awkwardly with autoregressive stepping; the coarse-level
approach above buys most of the global-coupling benefit without the per-step
iteration. Schwarz iteration remains a reasonable fallback if the coarse level
proves insufficient.

\section{Open questions for the next stage}

\begin{enumerate}[leftmargin=1.4em]
\item \textbf{Overlap vs.\ halo-forcing vs.\ consistent exchange} --- which merge
strategy gives acceptable seams at acceptable cost for autoregressive 3D rollout?
\item \textbf{Global coupling} --- is a coarse/global level required for the
elliptic pressure, and if so, coarse network, multipole path, or Schwarz
iteration?
\item \textbf{Minimum patch and halo size} --- set by the Kolmogorov scale, the
network's receptive field, and the geometry length scales of the urban canopy.
\item \textbf{Interface error accumulation} --- does the pushforward trick suffice
for stability when errors enter both at interfaces and over the rollout, or is a
generative/diffusion formulation \cite{lienen2024turbdiff,kohl2023benchmarking}
warranted?
\item \textbf{Training distribution} --- patches sampled from full-domain
simulations, with halos, so the network sees realistic interface conditions
during training rather than only at inference.
\end{enumerate}

\bibliographystyle{plain}
\bibliography{references}

\end{document}
